\subsubsection{Continuación regional y ecuación de renovación}

El selector \(169\) determina el impulso de grado seis. La prolongación
posterior se organiza por la transferencia \(D_A\) y una normalización
unitaria de la línea determinantal. Estas reglas admiten una formulación
en series formales antes de sustituir \(z=\alpha\), con \(D_A\) ya
construido. Escribamos
\[
 \mathscr C(z)=169z^6+R(z),\qquad R(z)\in z^7\mathbb R[[z]].
\]
La concatenación de una prolongación aumenta el grado en una unidad y
multiplica su peso por \(D_A\). La primera prolongación estricta tiene
peso \(D_A\). Por tanto, su ecuación de renovación es
\begin{equation}
 R(z)=D_Az^7+D_AzR(z).
 \label{eq:revision-renovacion}
\end{equation}

\begin{proposition}[Unicidad formal de la continuación normalizada]
La ecuación \eqref{eq:revision-renovacion} determina una única serie y
\[
 \mathscr C(z)=169z^6+\frac{D_Az^7}{1-D_Az}
             =169z^6+\sum_{n=7}^{\infty}D_A^{n-6}z^n.
\]
La realización es analítica en \(|D_Az|<1\).
\end{proposition}
\begin{proof}
El elemento \(1-D_Az\) tiene término independiente uno y es invertible
en el anillo de series formales. Resolver la ecuación produce la serie
exhibida. Equivalentemente, el coeficiente de grado siete es \(D_A\)
y cada coeficiente posterior es \(D_A\) veces el anterior. La serie
geométrica converge absolutamente en el disco indicado.
\end{proof}

La normalización de la primera prolongación es un dato operativo
indispensable. Si se conservaran únicamente el impulso de grado seis y
la razón de concatenación, las series
\[
 169z^6+\lambda\frac{D_Az^7}{1-D_Az}
\]
seguirían compartiendo ambos datos para cualquier \(\lambda\).
La regla unitaria fija \(\lambda=1\) antes de efectuar la evaluación.
Así queda localizada la contribución de cada regla a la unicidad del
lector, y la expresión racional conserva la continuación completa.

\subsubsection{Constante de Planck reducida y retorno de la sección de acción}

En la carta dimensional de \eqref{eq:el-secciones-accion}, la sección
anterior al retorno es la construcción del modelo destinada al contraste
con la constante de Planck reducida:
\[
 \hbar_{\mathrm{pre}}=H_5\,10^{-34}\mathcal U_S.
\]
La sección posterior conserva, además, la compensación regional
\(90\mathscr C(\alpha)/\pi\). Ambas pertenecen a la misma recta de
acción y su cociente \(R_{\mathrm{act}}\) mide el efecto multiplicativo
del retorno. El paso a \(h\) mediante \(2\pi\) utiliza la relación
entre la parametrización angular y una vuelta completa.

Esta construcción interviene efectivamente en
\(R_{\mathrm{act}}^{80-54\alpha+6\Delta_4}\). El electrón conserva
su coordenada central y su trayectoria; la escala cronológica no
reemplaza esa estructura por un cuanto indiferenciado. Tampoco se
multiplica de nuevo por \(10^{-34}\) una masa ya expresada en MeV:
la década pertenece a la sección de acción, cancela en su cociente
adimensional y reaparece únicamente donde actúa la realización
dimensional correspondiente.

El contraste numérico de \(H_5\), de la sección retornada y de
\(h/(2\pi)\) se presenta en la sección~\ref{sec:revision-planck-contraste}.
Así se distingue la denominación del objeto físico, la ecuación que el
modelo le asigna y la precisión con la que dicha ecuación lo reproduce.
